\documentclass[11pt,a4paper]{article}
\usepackage[margin=2.2cm]{geometry}
\usepackage{booktabs,multirow,longtable,graphicx,amsmath,hyperref,enumitem,listings,xcolor}
\usepackage[T1]{fontenc}
\usepackage{lmodern}
\hypersetup{colorlinks=true,linkcolor=blue!60!black,urlcolor=blue!60!black}
\setlength{\parskip}{4pt}\setlength{\parindent}{0pt}
\lstset{language=Python,basicstyle=\ttfamily\small,keywordstyle=\color{blue!70!black},commentstyle=\color{gray},
        stringstyle=\color{green!40!black},breaklines=true,frame=single,columns=fullflexible,showstringspaces=false}
\newcommand{\tbl}[1]{\IfFileExists{tables/#1.tex}{\input{tables/#1.tex}}{\emph{(table #1 appears after the run)}}}

\title{Study notes: Inverted Indexing, Sparse Retrieval and Vocabulary Mismatch\\
\large Everything behind the five parts of the assignment and every script in \texttt{cs6101/}}
\date{September 2026}

\begin{document}
\maketitle
\tableofcontents

\section{The whole pipeline in one picture}
Every part of the assignment is one stage of the same pipeline:
\begin{enumerate}[nosep]
  \item \textbf{Data}: BEIR corpora (documents), test queries, and \emph{qrels} (which documents are relevant to which query).
  \item \textbf{Index} (Part~1): turn the documents into a Lucene inverted index so a query can be answered in milliseconds.
  \item \textbf{Retrieve} (Part~2): score documents against a query with BM25 or TF-IDF and evaluate the ranking with nDCG, Recall, MRR, MAP.
  \item \textbf{Diagnose} (Part~3): look at the queries BM25 gets wrong and measure how much they share words with their gold document.
  \item \textbf{Expand the query} (Part~4): add words to the query, taken either from the top retrieved documents (Rocchio, RM3) or from text an LLM writes (HyDE).
  \item \textbf{Learn the expansion} (Part~5): SPLADE, a neural model that produces the sparse word-weight vector itself.
\end{enumerate}
All scripts live in one folder, share \texttt{common.py} and \texttt{datasets.py}, write JSON to \texttt{results/}, and
\texttt{make\_tables.py} turns the JSON into the LaTeX tables used by both reports. \texttt{run\_all.sh} runs the scripts in order,
one at a time (the laptop has 7.6\,GB of RAM for WSL and no GPU; two heavy jobs at once thrash or get killed by the OOM killer).

\section{Setup and how to run}
\begin{itemize}[nosep]
  \item WSL2 Ubuntu, Python 3.12, Java 21, \texttt{pyserini==2.3.0} (bundles Anserini and Lucene~10), \texttt{ir\_datasets}, \texttt{ir\_measures}, PyTorch (CPU), Transformers, matplotlib, scipy.
  \item \texttt{python download.py} once: fetches the SPLADE checkpoint, the Qwen model, and the three prebuilt SPLADE indexes.
  \item \texttt{bash run\_all.sh}: \texttt{index.py} $\to$ \texttt{index\_stats.py} $\to$ \texttt{phase2.py} $\to$ \texttt{phase3.py} $\to$ \texttt{phase4a.py} $\to$ \texttt{hyde\_generate.py} $\to$ \texttt{phase4b.py} $\to$ \texttt{phase5.py} $\to$ \texttt{make\_tables.py}. Log in \texttt{logs/run\_all.log}.
  \item Any script also takes dataset names: \texttt{python phase2.py scifact}.
  \item Compile: \texttt{pdflatex report.tex} (twice, for references) inside \texttt{report/}; same for \texttt{learning.tex}.
\end{itemize}
Two pitfalls met on this machine: a NUL byte written into the 1.7\,GB HotpotQA JSONL on the OneDrive mount made Anserini stop
silently at 2.9M of 5.2M documents (fix: delete and regenerate the file; check \texttt{Total N documents indexed} in the log),
and running the LLM, the SPLADE encoder and an index build together made the index build four times slower and triggered the
OOM killer. Timings in the report were therefore measured with one job at a time.

\section{Shared code: \texttt{datasets.py} and \texttt{common.py}}
\texttt{datasets.py} maps short names to \texttt{ir\_datasets} ids (\texttt{beir/scifact/test}, \ldots) and exposes
\texttt{docs\_iter()}, \texttt{queries\_iter()}, \texttt{qrels\_iter()}. A document has \texttt{doc\_id, title, text};
a query has \texttt{query\_id, text}; a qrel has \texttt{query\_id, doc\_id, relevance}.

\texttt{common.py}:
\begin{lstlisting}
METRICS = [nDCG @ 10, R @ 100, RR @ 10, AP]      # ir_measures objects
SAMPLE = 500                                       # queries for Parts 4-5 on FEVER / HotpotQA

def load(name, sample=None):        # queries + qrels, optionally a fixed random sample
    ...
    random.Random(0).shuffle(queries)               # seed 0 => same sample in every script
    queries = sorted(queries[:sample], key=lambda q: q.query_id)
def evaluate(run, qrels):           # run = {qid: {docid: score}}
    return {str(m): round(float(v), 4) for m, v in calc_aggregate(METRICS, qrels, run).items()}
def search(searcher, queries, k=100):
    return hits_to_run(searcher.batch_search([q.text ...], [q.query_id ...], k=k, threads=8))
\end{lstlisting}
Why a sample? Parts~4 and~5 do work \emph{per query} (a second search with a long weighted query, LLM generation, SPLADE
encoding). On 5M-document corpora that is minutes per hundred queries on a CPU, so a fixed 500-query subset keeps the
experiments feasible; the BM25 baseline is always recomputed on the same subset so comparisons stay fair.
\texttt{save(part, name, data)} merges one dataset's result into \texttt{results/<part>.json}.

\section{Part 1: building the index (\texttt{index.py}, \texttt{index\_stats.py})}
\subsection{Concept}
An \emph{inverted index} maps every term to a \emph{postings list}: the documents containing it and the term frequency (tf) in
each. Lucene also stores per document the length (for normalisation) and, only if asked (\texttt{-storeDocvectors}), the full
term vector of each document, which is what Rocchio needs later. Before indexing, text goes through an \emph{analyzer}:
tokenise, lowercase, remove stop words, Porter-stem (\emph{studies}, \emph{studying} $\to$ \emph{studi}). Queries go
through the same analyzer, otherwise nothing would match.

\subsection{Code}
\begin{lstlisting}
for doc in datasets.get_corpus(datasets.load_dataset(name)):
    f.write(json.dumps({"id": doc.doc_id, "contents": doc.title + "\n" + doc.text}) + "\n")
subprocess.run(["python", "-m", "pyserini.index.lucene", "-collection", "JsonCollection",
                "-input", collection_dir, "-index", index_dir, "-storeDocvectors"], check=True)
\end{lstlisting}
Pyserini's \texttt{JsonCollection} wants one JSON object per line with \texttt{id} and \texttt{contents}. The indexer is a
Java program, so it is launched as a subprocess and timed around the call; the directory size afterwards is the index size.
\texttt{index\_stats.py} then opens the index with \texttt{LuceneIndexReader}:
\begin{lstlisting}
reader.stats()                          # documents, unique_terms, total_terms
reader.get_document_vector(docid)       # {"cell": 5, "stem": 5, ...}  analysed terms -> tf
sum(vec.values())                       # document length in indexed tokens
reader.get_term_counts("cell", analyzer=None)[0]   # df: documents containing the term
\end{lstlisting}
\texttt{analyzer=None} matters: the term from a document vector is already stemmed; running the analyzer again could stem it a
second time and look up the wrong string.
\tbl{part1}\tbl{part1_demo}

\section{Part 2: BM25, tuning, TF-IDF (\texttt{phase2.py})}
\subsection{Concept}
BM25 scores a document $d$ for query $q$ as
\[ \sum_{t\in q} \mathrm{idf}(t)\,\frac{\mathrm{tf}(t,d)\,(k_1+1)}{\mathrm{tf}(t,d)+k_1\bigl(1-b+b\,|d|/\mathrm{avgdl}\bigr)} ,\qquad
   \mathrm{idf}(t)=\log\frac{N-\mathrm{df}(t)+0.5}{\mathrm{df}(t)+0.5}. \]
$k_1$ controls how quickly repeated occurrences stop adding score (saturation); $b\in[0,1]$ controls how strongly long documents
are penalised ($b{=}0$: no length normalisation, $b{=}1$: full). Pyserini's default is $k_1{=}0.9$, $b{=}0.4$.
Classic TF-IDF (Lucene's \texttt{ClassicSimilarity}) uses $\sqrt{\mathrm{tf}}\cdot\mathrm{idf}^2$ with a cosine-style length norm:
no saturation and a weaker, fixed length normalisation, which is why it is usually worse than BM25.

Metrics on a ranked list with relevance labels: \textbf{nDCG@10} = DCG of the top-10 ($\sum \mathrm{rel}_i/\log_2(i+1)$)
divided by the DCG of the ideal ordering; \textbf{Recall@100} = fraction of relevant documents found in the top-100;
\textbf{MRR@10} = $1/\mathrm{rank}$ of the first relevant document if it is within 10, else 0; \textbf{MAP} = mean over queries
of the average of precision at each relevant document's rank. All averaged over queries by \texttt{ir\_measures}.

\subsection{Code}
\begin{lstlisting}
searcher = LuceneSearcher(f"indexes/{name}")
searcher.set_bm25(0.9, 0.4)
for k1 in [0.6, 0.9, 1.2, 1.5, 2.0]:
    for b in [0.3, 0.4, 0.5, 0.75, 1.0]:
        searcher.set_bm25(k1, b)
        grid[(k1, b)] = common.evaluate(common.search(searcher, queries), qrels)
k1, b = max(grid, key=lambda p: grid[p]["nDCG@10"])
searcher.object.searcher.setSimilarity(ClassicSimilarity())     # Lucene IndexSearcher via pyjnius
\end{lstlisting}
The similarity is a property of the searcher, not of the index, so TF-IDF needs no re-indexing.
\tbl{part2}\tbl{part2_grid}

\section{Part 3: seeing vocabulary mismatch (\texttt{phase3.py})}
\subsection{Concept}
\emph{Vocabulary mismatch}: the query and the relevant document say the same thing with different words (synonyms:
\emph{tumour}/\emph{neoplasm}; paraphrase: \emph{reduces the risk of}/\emph{is protective against}; abbreviation:
\emph{MRI}/\emph{magnetic resonance imaging}). A lexical model such as BM25 can only match identical (stemmed) strings.
To measure it we use the \emph{Jaccard coefficient} $J=|Q\cap D|/|Q\cup D|$ between the set of analysed query tokens and
the set of terms in the gold document's indexed vector. Because both sides are analysed the same way, $J$ measures exactly
what BM25 could match.

How strongly does low $J$ predict a miss? We compare the $J$ distribution of pairs where the gold document is in the top-10
(success) with pairs where it is not (failure), and report the \emph{AUC}: the probability that a random success pair has a
higher $J$ than a random failure pair (0.5 = no signal, 1.0 = perfect). It equals the Mann--Whitney $U$ statistic divided by
$n_{\text{succ}}\,n_{\text{fail}}$, which is what \texttt{scipy.stats.mannwhitneyu} gives.

\subsection{Code}
\begin{lstlisting}
ranked = sorted(run[q.query_id], key=run[q.query_id].get, reverse=True)
rank = ranked.index(docid) + 1 if docid in ranked else None
success = rank is not None and rank <= 10
jaccard = len(q_tok & d_tok) / len(q_tok | d_tok)
auc = mannwhitneyu(succ, fail).statistic / (len(succ) * len(fail))
\end{lstlisting}
Failure categories are a simple heuristic: no shared term $\to$ \emph{no lexical overlap}; an all-caps token on one side
(regex \verb|\b[A-Z][A-Z0-9\-]{1,7}s?\b|) whose letters appear as the initials of consecutive words on the other side $\to$
\emph{abbreviation vs.\ expansion}; less than half of the query terms present $\to$ \emph{partial overlap}; otherwise
\emph{high overlap but outranked} (the words are there, other documents just score higher: length, competing near-duplicates,
or a multi-hop question whose second gold passage shares few words). The document text for the examples comes from
\texttt{ir\_datasets}' \texttt{docs\_store().get(docid)}, because the index does not store raw text.
\tbl{part3_stats}\tbl{part3_bins}\tbl{part3_cats}

\section{Part 4a: Rocchio and RM3 (\texttt{rocchio.py}, \texttt{phase4a.py})}
\subsection{Concept}
Pseudo-relevance feedback assumes the top-$N$ documents of a first search are relevant and uses their vocabulary to
re-write the query. Rocchio:
\[ w_t = \alpha\, f(q)[t] + \frac{\beta}{N}\sum_{d\in\text{top-}N}\tilde f(d)[t]. \]
$f(q)$ is the query's term-frequency vector, $\tilde f(d)$ the document's; we divide both by their length so that a long
document does not dominate and the query side keeps total weight $\alpha$. Terms present in more than 10\% of the corpus are
never added (they carry no information and would pull every query the same way); the $k$ heaviest remaining terms are kept.
\emph{Query drift}: when the top-$N$ contain off-topic documents, their words are added and the re-search moves away from the
original intent.

RM3 is the probabilistic cousin: it estimates a relevance language model from the feedback documents, keeps the top terms,
and interpolates it with the original query ($\lambda$ = \texttt{original\_query\_weight}). Pyserini implements it in Java
behind \texttt{set\_rm3}.

\subsection{Code}
\begin{lstlisting}
centroid = Counter()
for vec in vectors:                      # vectors: term -> tf, one per feedback document
    length = sum(vec.values()) or 1
    for t, tf in vec.items():
        centroid[t] += tf / length
weights = {t: alpha * tf / q_len + beta / n * centroid.get(t, 0.0) for t, tf in q.items()}
expansion = sorted((t for t in centroid if t not in q and not self.too_common(t)), key=lambda t: -centroid[t])[:k]
for t in expansion:
    weights[t] = beta / n * centroid[t]

builder = querybuilder.get_boolean_query_builder()
for term, w in weights.items():
    builder.add(querybuilder.get_boost_query(JTermQuery(JTerm("contents", term)), float(w)), JBooleanClauseOccur.should.value)
hits = searcher.search(builder.build(), k=100)
\end{lstlisting}
Why the query builder instead of pasting words into a string? A string query would be analysed again (double stemming) and
could not carry per-term weights; \texttt{BoostQuery} multiplies each term's BM25 contribution by $w_t$.
In \texttt{phase4a.py} the feedback vectors come from the index (\texttt{rocchio.index\_vectors(top-N docids)}), the run is
evaluated, per-query nDCG@10 is compared with the baseline to list drift candidates, and RM3 is run on the same queries:
\begin{lstlisting}
searcher.set_rm3(fb_terms=k, fb_docs=n_fb, original_query_weight=0.5)
entry["rm3"] = common.evaluate(common.search(searcher, queries), qrels)
searcher.unset_rm3()
\end{lstlisting}
\tbl{part4a}

\section{Part 4b: HyDE (\texttt{hyde\_generate.py}, \texttt{phase4b.py})}
\subsection{Concept}
HyDE (Hypothetical Document Embeddings) asks an LLM to write the document that \emph{would} answer the query, then uses that
text for retrieval. Here it replaces the corpus-retrieved feedback documents of 4a: the LLM's passage is analysed into a
term vector and fed to the same Rocchio code. Two ways to combine: \emph{naive concatenation} (query + passage as one bag of
words, so every generated word counts like a query word) versus the \emph{feedback model} (Rocchio keeps the query dominant,
filters common words and keeps only $k$ terms). The question the assignment asks is which matters more: where the words come
from (LLM vs.\ corpus) or how they are combined (concatenation vs.\ Rocchio).

\subsection{Code}
\begin{lstlisting}
tok = AutoTokenizer.from_pretrained(MODEL, padding_side="left")      # left padding for batched generation
model = AutoModelForCausalLM.from_pretrained(MODEL, dtype=torch.float32).eval()
prompts = [tok.apply_chat_template([{"role": "user", "content": PROMPT[name].format(q=q.text)}],
                                   tokenize=False, add_generation_prompt=True) for q in batch]
gen = model.generate(**inputs, max_new_tokens=max_new_tokens, do_sample=True, temperature=0.7, top_p=0.9,
                     num_return_sequences=n_docs, pad_token_id=tok.pad_token_id)
texts = tok.batch_decode(gen[:, inputs["input_ids"].shape[1]:], skip_special_tokens=True)   # drop the prompt tokens
\end{lstlisting}
The script appends one JSON line per query and skips ids already present, so it can be stopped and resumed. The model is
Qwen2.5-0.5B-Instruct because it is the largest that runs at a usable speed on this CPU (no GPU, so no vLLM).
In \texttt{phase4b.py}, \texttt{rocchio.text\_vectors(hyde[qid])} calls \texttt{Analyzer.compute\_document\_vector} on the
generated passages; everything after that is the unchanged 4a code.
\tbl{part4b}

\section{Part 5: SPLADE (\texttt{phase5.py})}
\subsection{Concept}
SPLADE puts a BERT encoder with its masked-language-model head over the text; for every input position $i$ the head produces
a logit $o_{ij}$ for every vocabulary word $j$. The representation is
\[ w_j = \max_i \log\bigl(1+\mathrm{ReLU}(o_{ij})\bigr), \]
one weight per vocabulary word, zero for most of them (the FLOPS regulariser during training pushes weights to zero). A word
gets weight even if it never appears in the text, because the MLM head predicts related words: that is learned expansion.
Documents are encoded once and stored in an \emph{impact index}: a Lucene index whose ``term frequency'' is the integer
weight (Pyserini quantises $w_j$ to integers), built with \texttt{-impact -pretokenized}. At query time the query is encoded
the same way and the score is the dot product, computed by the usual posting-list traversal. So SPLADE is neural on both
sides but still a sparse inverted index at retrieval time.

\subsection{Code}
\begin{lstlisting}
doc_encoder = SpladeDocumentEncoder(MODEL, device="cpu")
vectors = doc_encoder.encode([d["contents"] for d in batch], max_length=256)   # list of {word: int weight}
f.write(json.dumps({"id": d["id"], "contents": "", "vector": {t: w for t, w in v.items() if w > 0}}) + "\n")
subprocess.run(["python", "-m", "pyserini.index.lucene", "-collection", "JsonVectorCollection",
                "-input", coll, "-index", index, "-impact", "-pretokenized"], check=True)
searcher = LuceneImpactSearcher(index, SpladeQueryEncoder(MODEL, device="cpu"))
hits = searcher.batch_search(texts, qids, k=100, threads=8)
\end{lstlisting}
SciFact (5\,183 documents) is encoded locally; FEVER and HotpotQA would take days on a CPU, so their prebuilt
\texttt{beir-v1.0.0-*.splade-pp-ed} indexes (same checkpoint) are downloaded and only the queries are encoded here.
The expansion-term comparison lists SPLADE's non-zero words that are not literal word-pieces of the query
(\texttt{encoder.tokenizer.tokenize(text)}), next to Rocchio's and HyDE's terms for the same query; all three are passed
through the index analyzer (stemming) before counting overlaps, since SPLADE works on unstemmed word-pieces.
\tbl{part5}\tbl{part5_overlap}

\section{Reading the results}
\begin{itemize}[nosep]
  \item \texttt{results/part1.json}: build time, size; \texttt{part1\_stats.json}: IndexReader statistics.
  \item \texttt{part2.json}: default/tuned/TF-IDF metrics, the full grid, the chosen $k_1,b$, BM25 latency.
  \item \texttt{part3.json}: failure rate, Jaccard mean/median for success vs.\ failure, AUC, failure rate per Jaccard bin, category counts, examples; \texttt{part3\_jaccard\_*.pdf}: histograms.
  \item \texttt{part4a.json}: BM25 vs.\ Rocchio vs.\ RM3 per $(N,k)$, drift candidates, expansion terms of the first 12 queries.
  \item \texttt{hyde\_*.jsonl}: generated passages; \texttt{part4b.json}: concatenation vs.\ HyDE+Rocchio vs.\ corpus PRF.
  \item \texttt{part5.json}: SPLADE metrics, own-index encode/index time and size, per-query expansion rows and average overlaps.
\end{itemize}
What to expect and why: BM25 tuning moves nDCG@10 by a point or two at most; TF-IDF is clearly below BM25. Most BM25 failures
have low Jaccard, and pairs with $J{=}0$ almost always fail, but many failures still share words (they are \emph{outranked},
not unmatched), so overlap predicts failure strongly but not perfectly. Corpus PRF helps recall more than precision and can
hurt SciFact, whose claims are short and specific (drift). HyDE with Rocchio weighting tends to beat naive concatenation:
the combination method matters at least as much as the source. SPLADE, trained on MS~MARCO, is the strongest system on all
three datasets because its expansion is learned rather than guessed from a few documents.

\section{Deliverable checklist}
\begin{tabular}{p{0.55\textwidth}p{0.4\textwidth}}\toprule
Problem statement asks for & Where it is \\ \midrule
Corpus size, queries, deviations & report, Datasets paragraph \\
Index build time and size & Table part1 (\texttt{index.py}) \\
Default / tuned BM25 / TF-IDF table, $k_1,b$ and justification & Tables part2, part2\_grid, Part~2 text \\
Failure categories with examples, Jaccard comparison, verdict & Tables part3\_*, figure, Part~3 verdict \\
Rocchio and RM3 at $\geq 2$ $(N,k)$, drift examples & Table part4a, drift list \\
Concatenation / HyDE-Rocchio / corpus PRF, verdict & Table part4b, Part~4b verdict \\
SPLADE metrics, 10+ query three-way term table, overlap, disagreements & Tables part5, part5\_overlap, part5\_terms \\
LLM usage and chat links & last section of the report \\ \bottomrule
\end{tabular}

\section{Questions to test yourself}
\begin{enumerate}[nosep]
  \item Which call would fail if the index were built without \texttt{-storeDocvectors}?
  \item Why is Jaccard computed on analysed tokens rather than raw words?
  \item What happens to the Rocchio weights if you remove the division by document length?
  \item Why can the same Rocchio class be used for 4a and 4b without any change?
  \item Why does SPLADE need \texttt{-pretokenized}, and what would the analyzer do to its word-pieces otherwise?
  \item Given one query, compute its BM25 score for a two-term document by hand with $k_1{=}0.9$, $b{=}0.4$.
  \item Change the failure threshold to top-20: which numbers in Part~3 move, and in which direction?
\end{enumerate}
\end{document}
